\section{Scope of ordinal analysis}
\label{sec:normalization-curry-howard-04}

The landscape of ordinal-informative consistency analyses is surveyed in
\cite{Buss1998}.

\paragraph{Cited theorem.}
Gentzen's consistency analysis of arithmetic, in the sequent-reduction
presentation surveyed in
\cite[Chapter~10, especially Sections~10.1--10.4, pp.~318--337]
     {TroelstraSchwichtenberg2000}, assigns notations below
\(\varepsilon_0\) to arithmetic derivations so that each critical reduction
strictly lowers the assigned notation.  In a metatheory that accepts
transfinite induction for the predicates and primitive-recursive ordering used
in that particular reduction analysis, termination yields
\(\mathsf{Con}(\PA)\).  This is an analysis-relative imported result;
it is not a definition of a canonical invariant and it does not assert that an
unspecified ordinal ``measures'' a theory.

\paragraph{Formal objects needed to read the citation.}
The symbol \(\varepsilon_0\) denotes the least nonzero ordinal
\(\alpha\) satisfying \(\omega^\alpha=\alpha\).  A standard external way to
write ordinals below it uses finite Cantor-normal-form terms
\[
  a=\omega^{a_1}\cdot m_1+\cdots+
    \omega^{a_k}\cdot m_k,
\]
where \(k\geq1\), each \(m_i\) is a positive natural number, the exponent
terms satisfy \(a_1>\cdots>a_k\), and \(0\) is also a term.  Exponents are
formed recursively by the same clause.  Interpreting the symbols by ordinal
addition and exponentiation assigns each such finite term an ordinal below
\(\varepsilon_0\).  This display supplies notation for the following worked
calculation; it is not asserted to be the source's complete primitive-recursive
coding.

For example, \(\lambda=\omega^\omega\) is a limit notation.  The sequence
\[
  \lambda[n]=\omega^n\qquad(n\geq1)
\]
satisfies \(\lambda[n]<\lambda\) for every \(n\) and
\(\sup_{n\geq1}\lambda[n]=\lambda\).  Likewise, for
\(\mu=\omega^3\), define \(\mu[n]=\omega^2\cdot n\) for \(n\geq1\).  Then
\[
  \mu[4]=\omega^2\cdot4<\omega^3=\mu.
\]
This is the elementary shape of a fundamental-sequence descent: a limit
notation is replaced by a selected smaller approximation.  It does not show
that any particular proof reduction receives \(\mu\), because that claim
depends on the source's derivation coding and ordinal assignment.

The cited result therefore has three separately typed ingredients: finite
derivation codes for the selected sequent presentation of first-order
arithmetic; ordinal-notation codes with a primitive-recursive comparison and
fundamental-sequence operation; and a primitive-recursive reduction on
derivation codes together with an assignment from derivations to notation
codes.  Its decisive lemma says, case by case for that reduction, that a
critical reduction strictly decreases the assigned notation.  Transfinite
induction for the predicates used in the reduction then excludes an infinite
reduction sequence, and the source extracts the consistency statement from
the resulting normal form.

\paragraph{Internal proof boundary.}
A complete proof in this manuscript would still have to reproduce the exact
arithmetic rule surface and derivation coding, the complete notation grammar
and algorithms, every reduction clause, the assignment, every case of the
decrease lemma, the precise transfinite-induction schema, and the final
consistency extraction.  A later use of the technical phrase
\emph{proof-theoretic ordinal} for an arbitrary theory would additionally need
a stated reduction notion, a leastness or supremum clause, and a comparison
showing what is invariant under alternative presentations.  No such general
invariant is defined here.

The imported result is relative.  It does not give a finitary proof of
\(\mathsf{Con}(\PA)\) formalizable in \(\PA\) itself.  If \(\PA\) is
consistent, the second incompleteness theorem excludes such an internal proof
for the standard arithmetized consistency statement.  Stronger theories can
require notation systems far beyond \(\varepsilon_0\), and some strength
comparisons are not captured by a single canonical ordinal.

\begin{warning}[Unqualified strength claims]
Do not infer that a larger displayed ordinal alone makes one theory stronger
than another.  Compare the theories, notation systems, base metatheories, and
reflection principles used in the analyses.
\end{warning}

